\section{Results}
\label{sec:results}

\paragraph{Aggregate test split.} On the mixed test split, which blends all five
query styles, MoC-RAG with $\mathrm{top\_experts}{=}3$ posts the best aggregate
Recall@8 (86\%) and the strongest MRR/nDCG, ahead of BM25 (78\%) and the dense
family (78--82\%), while carrying roughly half their hard distractors at
95--100\% routing accuracy (Table~\ref{tab:main}). We stress that this aggregate
already reflects the harder query mix; on a purely keyword-aligned split BM25 is
co-leading, as the robustness breakdown makes explicit.

\begin{table}[t]
  \centering
  \caption{Main benchmark results on the mixed test split (real embedder).
  MoC-RAG rows shaded. HardDistr and Tokens: lower is better.}
  \label{tab:main}
  \begin{tabular}{l rrrr rr rr}
\toprule
Method & Recall@8 & Prec@8 & MRR & nDCG@8 & HardDistr$\downarrow$ & Tokens$\downarrow$ & UsefulRatio$\uparrow$ & RouteAcc \\
\midrule
bm25\_rag & 78\% & 19\% & 0.719 & 0.727 & 167 & 23560 & 12\% & -- \\
dense\_rag & 78\% & 19\% & 0.558 & 0.617 & 271 & 23539 & 12\% & -- \\
hybrid\_rag & 82\% & 21\% & 0.723 & 0.740 & 244 & 23519 & 12\% & -- \\
metadata\_rag & 82\% & 20\% & 0.703 & 0.721 & 273 & 23364 & 13\% & -- \\
reranked\_rag & 82\% & 21\% & 0.717 & 0.745 & 249 & 23499 & 12\% & -- \\
\rowcolor{mocblue!8} moc\_rag\_e1 & 81\% & 20\% & 0.748 & 0.760 & 104 & 23509 & 12\% & 84\% \\
\rowcolor{mocblue!8} moc\_rag\_e2 & 85\% & 21\% & 0.768 & 0.781 & 122 & 23547 & 13\% & 95\% \\
\rowcolor{mocblue!8} moc\_rag\_e3 & 86\% & 22\% & 0.789 & 0.796 & 156 & 23529 & 13\% & 99\% \\
\rowcolor{mocblue!8} moc\_rag\_all & 85\% & 21\% & 0.766 & 0.779 & 196 & 23568 & 12\% & 100\% \\
\bottomrule
\end{tabular}

\end{table}

\paragraph{Robustness by query type.} Table~\ref{tab:robust} and
Figures~\ref{fig:recall_by_type}, \ref{fig:drop} report the central result.
BM25's Recall@8 falls from 100\% (keyword) to 64\% (adversarial), a 36-point
collapse, because misleading keywords pull it toward the contradictory negative.
MoC-RAG holds within $\sim$17 points (96\%$\to$79\% for $k{=}3$) and overtakes
BM25 on the adversarial split by $\sim$15 points. The dense, hybrid,
metadata-filtered, and reranked baselines are less brittle than BM25 on recall
but pay for it in harmful context.

\begin{table}[t]
  \centering
  \caption{Robustness by query type (Recall@8 and Hard Distractors per split;
  $\Delta$ is the keyword$\to$adversarial recall change in points).}
  \label{tab:robust}
  \begin{tabular}{l rrr rrr r}
\toprule
Method & R@8 key. & R@8 par. & R@8 adv. & HD key. & HD par. & HD adv. & $\Delta$ kw$\to$adv \\
\midrule
bm25\_rag & \textbf{100\%} & 81\% & 64\% & 27 & 68 & 72 & -36 \\
dense\_rag & 83\% & 79\% & 74\% & 64 & 104 & 103 & -9 \\
hybrid\_rag & 96\% & 88\% & 71\% & 52 & 101 & 91 & -25 \\
metadata\_rag & 92\% & 88\% & 71\% & 63 & 107 & 103 & -21 \\
reranked\_rag & 83\% & 86\% & 78\% & 56 & 96 & 97 & \textbf{-5} \\
\rowcolor{mocblue!8} moc\_rag\_e1 & \textbf{100\%} & 84\% & 69\% & \textbf{22} & \textbf{40} & \textbf{42} & -31 \\
\rowcolor{mocblue!8} moc\_rag\_e2 & 96\% & 86\% & 78\% & 26 & 48 & 48 & -18 \\
\rowcolor{mocblue!8} moc\_rag\_e3 & 96\% & \textbf{89\%} & \textbf{79\%} & 30 & 64 & 62 & -17 \\
\rowcolor{mocblue!8} moc\_rag\_all & 98\% & 85\% & 77\% & 34 & 82 & 80 & -21 \\
\bottomrule
\end{tabular}

\end{table}

\paragraph{Harmful context.} Figure~\ref{fig:hd} shows packed hard distractors on
the adversarial split: MoC-RAG carries 48--62 versus 91--103 for the dense
family---roughly half. Figure~\ref{fig:pareto} places methods on a recall vs.\
hard-distractor Pareto plane; MoC-RAG variants dominate the upper-right (more
recall, less harmful context).

\begin{figure}[t]
  \centering
  \begin{minipage}{0.49\linewidth}
    \centering\includegraphics[width=\linewidth]{adversarial_recall_drop.pdf}
    \caption{Recall change, keyword$\to$adversarial (closer to 0 is more robust).}
    \label{fig:drop}
  \end{minipage}\hfill
  \begin{minipage}{0.49\linewidth}
    \centering\includegraphics[width=\linewidth]{hard_distractors.pdf}
    \caption{Hard distractors packed on the adversarial split (lower is better).}
    \label{fig:hd}
  \end{minipage}
\end{figure}

\begin{figure}[t]
  \centering
  \begin{minipage}{0.49\linewidth}
    \centering\includegraphics[width=\linewidth]{recall_vs_distractors.pdf}
    \caption{Pareto front: adversarial recall vs.\ hard distractors.}
    \label{fig:pareto}
  \end{minipage}\hfill
  \begin{minipage}{0.49\linewidth}
    \centering\includegraphics[width=\linewidth]{routing_accuracy.pdf}
    \caption{Expert routing accuracy by MoC-RAG variant.}
    \label{fig:route}
  \end{minipage}
\end{figure}

\paragraph{Ablation over expert fan-out.} Table~\ref{tab:abl} varies
$\mathrm{top\_experts}$. Routing accuracy rises monotonically with $k$
(Figure~\ref{fig:route}), reaching 100\% at $k{=}3$ and \texttt{all}; $k{=}2$ and
$k{=}3$ give the best recall/robustness trade-off, while $k{=}1$ minimizes hard
distractors at some recall cost. Widening to all experts recovers a flat-RAG-like
candidate pool and erodes the distractor advantage---evidence that the
\emph{selection}, not merely the typed store, drives the gains.

\begin{table}[t]
  \centering
  \caption{Ablation: MoC-RAG expert fan-out $k$ across query types.}
  \label{tab:abl}
  \begin{tabular}{l rrr r r}
\toprule
Variant & R@8 key. & R@8 par. & R@8 adv. & HD adv. & RouteAcc \\
\midrule
$k{=}1$ & 100\% & 84\% & 69\% & 42 & 75\% \\
$k{=}2$ & 96\% & 86\% & 78\% & 48 & 96\% \\
$k{=}3$ & 96\% & 89\% & 79\% & 62 & 100\% \\
all & 98\% & 85\% & 77\% & 80 & 100\% \\
\bottomrule
\end{tabular}

\end{table}
